\scanpage{187}
\mspage{61}
% archive sheet number: 186

par $\underline{\lim}$ de cardinal $\leq\unsure{\pi|\Pi}$. Donc si $J$ est une partie $\pi$-cofinale de $\mathrm{Ss}_{\pi}(X)$, $J\cap J'$ est cofinal dans $J$, i.e. dans $I$, sss $J'$ \deleted{\unsure{$\in$}} est \deleted{car alors elle sera} cofinal dans $I$ (i.e. sss $X\in M$, par la condition c). Car $J\cap J'$ cofinal $\Rightarrow J'$ cofinal tautologiquement, et $J'$ cofinal implique $J'$ $\pi$-cofinal, donc $J\cap J'$ $\pi$-cofinal par le lemme 5.8, donc cofinale.

Reste à prouver seulement que b) $\Rightarrow$ a), i.e. la stabilité de $M$ par limites inductives \inserted{filtrantes} quelconques, moyennant \deleted{le fait} b). Ici à nouveau, cela ne doit pas utiliser l'hypothèse particulière sur $\pi$, ni même l'hypothèse d'exactitude e) sur $\mathcal M$, mais être une histoire générale d'accessibilité :

\uline{Proposition} 5.9. Soit $\mathcal M$ une \deleted{partie} catégorie satisfaisant a) b) c) d) (p\unsure{ges} 44,45), $\pi\geq\pi_0$ un cardinal, $M$ une partie $\pi$-accessible de $\mathrm{Ob}\,\mathcal M$, i.e satisfaisant la condition \unsure{$1.b_{\pi})$} de \uline{$\overline{\mathrm{XIII}}$} p88. Alors $M$ est stable par $\varinjlim$ \unsure{filtrantes|filtrante} sss $M^\pi\overset{\mathrm{def}}=M\cap\mathrm{Filt}^{\pi}$ est stable par $\varinjlim$ filtrantes de cardinal $\leq\pi$.

Je n'ai pas envie d'expliciter ici la démonstration, ayant l'impression pénible de refaire à \unsure{perpète} la même démonstration (ici et dans \uline{$\overline{\mathrm{XIII}}$}). Je \unsure{devrai|devrais} revenir là-dessus dans une \unsure{rédaction|situation} plus décantée.

